% acknowledgements.tex

本科毕业设计即将完成，回顾整个过程，从最开始理解 TimeMixer 模型原理，到完成模块改进、实验整理和论文写作，我对时序数据预测模型的理解比课题开始时更加具体，也更加认识到实验验证和结果分析的重要性。

首先感谢指导教师王鸽老师在毕业设计过程中给予的指导和帮助。老师在选题方向、研究思路和论文写作等方面提出了许多建议，使我能够逐步明确课题目标，并在遇到实验结果不稳定、论文口径需要调整等问题时及时修正方向。

感谢电子与信息学部和计算机科学与技术专业在本科阶段提供的课程训练和学习环境。相关课程中学习到的机器学习、深度学习和程序设计知识，为本课题的模型理解、结构修改和实验分析提供了基础。

感谢同学和朋友在毕业设计过程中给予的交流和帮助。在实验运行、结果整理和论文排版过程中，与同学的讨论帮助我发现了一些容易忽视的问题，也让我在持续修改中保持耐心。

最后感谢家人在学习和生活上的支持。毕业设计是本科阶段一次较完整的综合训练，虽然本文仍有不足，但这个过程让我对科研和工程实践都有了更真实的认识。今后我会继续保持认真学习和实事求是的态度，在后续学习和工作中不断提高自己的能力。
